
\section{拉格朗日嵌入的障碍}\label{sec:9-2}

以下把 $\mathbb C^n$ 中的向量写成 $z=(z_1,\ldots,z_n)$，并约定
$\mathbb C^n$ 配备标准\BookRecordIndex{T0363}辛形式
\[
  \omega_0:=\sum_{j=1}^n dx_j\wedge dy_j,
  \qquad z_j=:x_j+iy_j,
\]
又以
\[
  \mathbb D:=\{s+it\in\mathbb C\mid s^2+t^2\le 1\}
\]
表示 $\mathbb C$ 中的闭单位圆盘。Gromov 在 \srcite{160} 中证明了\BookRecordIndex{T0222}黎曼映射定理的如下
推广。

\begin{theorem}[Gromov]\label{thm:9-2-1}
设 $L\subset\mathbb C^n$ 是紧\BookRecordIndex{T0219}拉格朗日子流形。则存在非常值全纯圆盘
$u:\mathbb D\to\mathbb C^n$，使得 $u(\partial\mathbb D)\subset L$。
\end{theorem}

他还对每个 $n>1$ 构造了 $\mathbb C^n$ 上的一个\BookRecordIndex{T0363}辛形式 $\omega$，它容许一个如下意义下
的闭\BookRecordIndex{T0401}正合拉格朗日子流形 $L$：$\omega$ 的某个原始 $1$-形式 $\lambda$ 在 $L$ 上的限制
$\lambda|_L$ 是正合的。（由于 $1$-形式 $\lambda$ 满足 $d\lambda=\omega$，这个限制当然
是闭的。）定理~\ref{thm:9-2-1} 蕴含这一辛结构 $\omega$ 是奇异的，即
$(\mathbb C^n,\omega)$ 不能嵌入标准欧氏空间 $(\mathbb C^n,\omega_0)$。关于这个主题的
进一步讨论见 \srcite{277}，近期的相关结果见 McLean \srcite{283}。

该定理还蕴含：欧氏空间 $(\mathbb C^n,\omega_0)$ 中的任意嵌入\BookRecordIndex{T0219}拉格朗日子流形 $L$ 都有
$H^1(L;\mathbb R)\ne0$。事实上，$1$-形式
\[
  \lambda:=\frac12\sum_{j=1}^n(x_j\,dy_j-y_j\,dx_j)
\]
在 $L$ 上的限制沿环路
$\mathbb R/2\pi\mathbb Z\to L:\theta\mapsto u(e^{i\theta})$ 的积分非零，因而代表一个
非零上同调类。相比之下，确有单连通的浸入拉格朗日流形。例如浸入
\[
  \iota:S^n\longrightarrow\mathbb C^n,
  \qquad
  \iota(x_0,x_1,\ldots,x_n)
  :=(x_1,x_0x_1,\ldots,x_n,x_0x_n)
\]
就是拉格朗日的，并把两点 $(\pm1,0,\ldots,0)$ 映到原点。

不过，关于哪些流形能拉格朗日嵌入欧氏空间，定理~\ref{thm:9-2-1} 并未提供太多其他
信息，这个问题至今仍未完全解决。可以作出一些简单观察。例如，若
$L\subset\mathbb C^n$ 是任意定向嵌入\BookRecordIndex{T0219}拉格朗日子流形，则其 Euler 示性数必为零。为此
注意，\BookRecordIndex{T0117}法丛通过
\[
  T_zL\longrightarrow T_zL^\perp:\zeta\longmapsto i\zeta
\]
同构于切丛。因此 $L$ 的自交数等于
\[
  L\mathbin{\boldsymbol\cdot}L=(-1)^{n(n-1)/2}\chi(L).
\]
但 $\mathbb C^n$ 的每个闭可定向子流形的自交数均为零。因此，环面是仅有的能拉格朗日
嵌入 $\mathbb C^2$ 的可定向曲面。在不可定向情形，同一论证表明 Euler 示性数必须为
偶数。具有偶 Euler 示性数的不可定向曲面是 $g+1$ 个 Klein 瓶的\BookRecordIndex{T0227}连通和：
\[
  K_g:=K\mathbin{\#}\cdots\mathbin{\#}K
  \cong K\mathbin{\#}\Sigma_g,
  \qquad \chi(K_g)=-2g.
\]
人们早已知道，对每个 $g>0$，$K_g$ 都能拉格朗日嵌入 $\mathbb C^2$；见 Givental
\srcite{144}，亦见 Lalonde--Sikorav \srcite{230}。$g=0$ 的情形最近由 Nemirovski
\srcite{298} 解决：他证明 Klein 瓶不能拉格朗日嵌入 $\mathbb C^2$。

\begin{remark}\label{rem:9-2-2}
此后，定理~\ref{thm:9-2-1} 得到了细化：人们找到了边界落在 $L$ 上的全纯圆盘之\BookRecordIndex{T0242}马斯洛夫
数的估计。Audin 在 \srcite{22} 中猜想，环面的每个拉格朗日嵌入
$\iota:\mathbb T^n\to\mathbb C^n$ 都容许一个\BookRecordIndex{T0242}马斯洛夫数为二的（不必全纯的）圆盘
\[
  u:(\mathbb D,\partial\mathbb D)\longrightarrow(\mathbb C^n,L),
  \qquad L:=\iota(\mathbb T^n).
\]
Viterbo \srcite{403} 对 $n=2$ 证明了这一点。对单调\BookRecordIndex{T0219}拉格朗日子流形
$L\subset\mathbb C^n$，Polterovich \srcite{324} 证明存在\BookRecordIndex{T0242}马斯洛夫数
$\mu\le n+1$ 的全纯圆盘；Oh \srcite{304} 在 $n\le24$ 的假设下找到\BookRecordIndex{T0242}马斯洛夫数为 $2$
的全纯圆盘（仍在单调情形）；Seidel 又在一项未发表工作中把后一结果推广到所有维数。
Biran 和 Cieliebak 在 \srcite{42} 中证明，对 $\mathbb C^n$ 的一般紧\BookRecordIndex{T0219}拉格朗日子流形，乃至
次临界 Stein 流形中的一般紧\BookRecordIndex{T0219}拉格朗日子流形，都存在\BookRecordIndex{T0242}马斯洛夫数满足
\[
  3-n\le\mu\le n+1
\]
的全纯圆盘。（关于不等式 $\mu\le n+1$，亦见下文注~\ref{rem:9-2-5}。）Buhovsky
\srcite{50} 已经对 $\mathbb R^{2n}$ 中的单调环面证明了 Audin 猜想。更近一些，Damian
\srcite{77} 把 Audin 猜想推广到 $\mathbb C^n$ 中一类更一般的单调\BookRecordIndex{T0219}拉格朗日子流形。不过，
截至本书写作时，非单调的 Audin 猜想仍未解决。另一个相关的公开问题（由 Paul Biran
向作者指出）是：$\mathbb C^n$ 的每个\BookRecordIndex{T0219}拉格朗日子流形，不论是否为环面，其\BookRecordIndex{T0241}马斯洛夫类是否
都非零。
\end{remark}


定理~\ref{thm:9-2-1} 的证明思路，是研究偶对 $(\lambda,u)$ 所满足的如下扰动
\BookRecordIndex{T0209}柯西--黎曼方程，其中 $0\le\lambda\le1$，而
$u:\mathbb D\to\mathbb C^n$ 是边界落在 $L$ 上的圆盘：
\begin{equation}\label{eq:9-2-1}
  \partial_su+J_0\partial_tu
  =\nabla H^\lambda_{s,t}(u),
  \qquad u(\partial\mathbb D)\subset L,
  \qquad [u]=A_0.
\end{equation}
这里 $A_0\in\pi_2(\mathbb C^n,L)$ 表示常值映射的同伦类，$\nabla$ 表示相对于标准度量的
梯度。函数族 $H^\lambda$（$\lambda\in[0,1]$）由
\[
  [0,1]\times\mathbb D\times\mathbb C^n\longrightarrow\mathbb R,
  \qquad(\lambda,s,t,z)\longmapsto H^\lambda_{s,t}(z)
\]
给出；它从 $H^0_{s,t}(z)=0$ 出发，以线性映射
$H^1_{s,t}(z)=\langle a,z\rangle$ 结束，其中 $a\in\mathbb C^n$ 是适当选取的向量；并且
存在常数 $c>0$，使所有中间的 $\lambda$ 都满足
\begin{equation}\label{eq:9-2-2}
  H^\lambda_{s,t}(z)=\lambda\langle a,z\rangle,
  \qquad |z|\ge c.
\end{equation}
我们将在引理~\ref{lem:9-2-4} 中证明，适当选择 $a,c$ 后，方程~\eqref{eq:9-2-1} 在
$\lambda=1$ 时无解。另一方面，当 $\lambda=0$ 时，通过 $L$ 中每一点恰有一个解（常值
映射）。我们还会看到，对一般的 $H$，方程~\eqref{eq:9-2-1} 的全部解
$(\lambda,u)$（$\lambda\in[0,1]$）中通过 $L$ 内一个一般点的那些构成一维流形。关键
在于：这个解空间不可能紧，因为若它紧，它就是只有一个边界点的紧一维流形。但我们会
看到，非\BookRecordIndex{T0192}紧性只能来自 $J$-全纯圆盘或球面的\BookRecordIndex{T0243}冒泡。由于 $\mathbb C^n$ 不含 $J$-全纯球面，
必然存在边界落在 $L$ 上的某个 $J$-全纯圆盘。

\begin{remark}\label{rem:9-2-3}
注意，这里考察的是满足拉格朗日边界条件的 $J$-全纯圆盘。虽然第~4 章讨论过这类圆盘的
\BookRecordIndex{T0243}冒泡，但第~3 章没有明确处理相应的边值问题。不过，如下文所示，它和闭曲线情形没有太大
区别。在我们的应用中，主要问题是验证线性化方程是 Fredholm 的；这由定理 C.1.10
给出。在当前情形，像第~8 章那样作\BookRecordIndex{T0157}哈密顿扰动，便足以实现
\BookRecordIndex{T0166}横截性。（另一种做法是像第~\ref{sec:6-7} 节那样让 $J$ 依赖于定义域中的点。）若所处的
几何情形要求在 $M$ 上找到一般的\BookRecordIndex{T0088}殆复结构，那么正如第~\ref{sec:3-2} 节所做，需要考察
某处单射圆盘。然而，无处单射圆盘的结构远比第~\ref{sec:2-5} 节所述闭曲线情形复杂，
因为圆盘会以错综复杂的方式缠绕自身。Lazzarini \srcite{234,235} 和 Kwon--Oh
\srcite{220} 的论文解释了一般圆盘的结构（亦见 Zehmisch \srcite{425}）；
Biran--Cornea \srcite{43, §3} 利用这一结构找到了这样的一般 $J$。
\end{remark}

下面给出定理~\ref{thm:9-2-1} 证明中的第一个关键引理。

\begin{lemma}\label{lem:9-2-4}
设 $a\in\mathbb C^n$，固定常数 $c\ge\sup_{z\in L}|z|$，并设
$u:\mathbb D\to\mathbb C^n$ 满足 $u(\partial\mathbb D)\subset L$。

\begin{enumerate}[label=\textup{(\roman*)}]
\item 若 $u$ 满足方程 $\partial_su+J_0\partial_tu=a$，则 $|a|\le2c$。
\item 若 $\mathbb D\times\mathbb C^n\to\mathbb R:(s,t,z)\mapsto H_{s,t}(z)$ 是光滑
函数，满足
\[
  |z|\ge c\quad\Longrightarrow\quad H_{s,t}(z)=\langle a,z\rangle,
\]
并且 $u$ 满足方程
$\partial_su+J_0\partial_tu=\nabla H_{s,t}(u)$，则
$\sup_{s,t}|u(s,t)|\le c$。
\end{enumerate}
\end{lemma}

\begin{proof}
第二个断言的证明基于如下观察：方程
$\partial_su+J_0\partial_tu=a$ 的每个解都是调和的，因而满足
\[
  \Delta|u|^2=2|\partial_su|^2+2|\partial_tu|^2\ge0,
\]
其中 $\Delta:=\partial_s^2+\partial_t^2$ 表示 Laplace 算子。这说明 $|u|^2$ 在开集
\[
  \Omega:=\{s+it\in\mathbb D\mid |u(s,t)|>c\}
\]
上是次调和的。由于 $c\ge\sup_{z\in L}|z|$，该集合包含于 $\mathbb D$ 的内部。我们断言
$\Omega=\varnothing$。否则，由\BookRecordIndex{T0079}次调和函数的\BookRecordIndex{T0266}平均值不等式，$|u|$ 在 $\Omega$ 的每个连通
分支上都恒定；但在 $\partial\Omega$ 上 $|u|=c$，这是不可能的。这证明了 (ii)。为证明
(i)，由散度定理，
\[
\begin{aligned}
  |a|
  &=\frac1\pi\left|\int_{\mathbb D}
      (\partial_su+J_0\partial_tu)\,ds\,dt\right| \\
  &=\frac1\pi\left|\int_0^{2\pi}
      (\cos\theta+\sin\theta J_0)u(e^{i\theta})\,d\theta\right| \\
  &\le\frac1\pi\int_0^{2\pi}|u(e^{i\theta})|\,d\theta
  \le2c.
\end{aligned}
\]
引理得证。
\end{proof}

\begin{proof}[定理~\ref{thm:9-2-1} 的证明]
反设不存在非常值全纯圆盘
$u:(\mathbb D,\partial\mathbb D)\to(\mathbb C^n,L)$。任取 $a\in\mathbb C^n$ 使
\[
  |a|>c,
  \qquad c:=2\sup_{z\in L}|z|.
\]
记 $\mathcal H\subset C^\infty([0,1]\times\mathbb D\times\mathbb C^n)$ 为满足
\[
  H^0_{s,t}(z)=0,
  \qquad H^1_{s,t}(z)=\langle a,z\rangle,
\]
以及式~\eqref{eq:9-2-2} 的光滑哈密顿函数 $H$ 之集合。对每个
$H\in\mathcal H$，考察\BookRecordIndex{T0246}模空间
\[
  \mathcal M(H):=\bigl\{(\lambda,u)\bigm|0\le\lambda\le1,
  \ u:\mathbb D\to\mathbb C^n\text{ 满足式~\eqref{eq:9-2-1}}\bigr\}.
\]
我们证明，对每个 $H\in\mathcal H$，$\mathcal M(H)$ 都是紧的。$\mathbb C^n$ 中没有
全纯球面；依假设，其中也没有边界落在 $L$ 上的全纯圆盘。此外，由引理
~\ref{lem:9-2-4}(ii)，方程~\eqref{eq:9-2-1} 的每个解
$u:\mathbb D\to\mathbb C^n$ 都取值于紧集
$\{z\in\mathbb C^n\mid |z|\le c\}$。\BookRecordIndex{T0249}能量恒等式为
\[
\begin{aligned}
  E(u)&:=\int_{\mathbb D}|\partial_su|^2\,ds\,dt\\
  &=\int_{\mathbb D}
    \bigl\langle\partial_su,-J_0\partial_tu
      +\nabla H^\lambda_{s,t}(u)\bigr\rangle\,ds\,dt\\
  &=\int_{\mathbb D}\omega_0(\partial_su,\partial_tu)\,ds\,dt
    +\int_{\mathbb D}\partial_s(H^\lambda_{s,t}\circ u)\,ds\,dt
    -\int_{\mathbb D}(\partial_sH^\lambda_{s,t})\circ u\,ds\,dt\\
  &=\int_0^{2\pi}\cos\theta\,
      H^\lambda_{e^{i\theta}}(u(e^{i\theta}))\,d\theta
    -\int_{\mathbb D}(\partial_sH^\lambda_{s,t})\circ u\,ds\,dt,
\end{aligned}
\]
其中最后一个等式用到了 $[u]=0$。因此，\BookRecordIndex{T0246}模空间 $\mathcal M(H)$ 的元素有一致能量上界。
现在把定理~\ref{thm:4-6-1} 应用于紧\BookRecordIndex{T0358}辛流形
\[
  \widetilde M:=\mathbb D\times\{z\in\mathbb C^n\mid |z|\le c\},
\]
由 $J_0,H$ 经式~\eqref{eq:8-1-6} 诱导的 $\widetilde M$ 上\BookRecordIndex{T0088}殆复结构
$\widetilde J$，以及\BookRecordIndex{T0219}拉格朗日子流形 $\widetilde L:=S^1\times L$，即可\BookRecordIndex{T0326}推出
$\mathcal M(H)$ 紧。这里用到如下事实：正如第~\ref{sec:6-7} 节与第~\ref{sec:8-5} 节，
若定义域的一点 $z$ 处发生\BookRecordIndex{T0243}冒泡，则\BookRecordIndex{T0275}气泡的像必位于 $z$ 上方的纤维中，因而必为全纯圆盘
或球面。

其次注意，定理~\ref{thm:8-3-1} 与定理~\ref{thm:8-3-3} 的证明可以直接移植到当前
情形，并说明：对一般的\BookRecordIndex{T0157}哈密顿扰动 $H\in\mathcal H_{\mathrm{reg}}$，\BookRecordIndex{T0246}模空间
$\mathcal M(H)$ 是光滑的 $n+1$ 维流形。关键是，方程~\eqref{eq:9-2-1} 的每个解
$u:\mathbb D\to\mathbb C^n$ 必通过区域 $\{|z|<c\}$，我们可以在其中扰动 $H$；这是因为
$u(\partial\mathbb D)\subset L$，而 $L$ 包含在该区域中。明确这一点后，通常的论证即
适用：先证明由三元组 $(\lambda,u,H)$ 组成的\BookRecordIndex{T0268}普适模空间
$\mathcal M(\mathcal H^\ell)$ 是 $C^{\ell-1}$ Banach 流形，其中
$H\in\mathcal H^\ell$ 且 $(\lambda,u)\in\mathcal M(H)$；再考察投影
\[
  \pi:\mathcal M(\mathcal H^\ell)\longrightarrow\mathcal H^\ell
\]
的\BookRecordIndex{T0409}正则值集合 $\mathcal H^\ell_{\mathrm{reg}}$。上述\BookRecordIndex{T0192}紧性结果表明该投影是恰当的，因此其
\BookRecordIndex{T0409}正则值集合是开集；而采用 萨德--斯梅尔定理的通常论证表明该集合稠密。进一步，由定理
C.1.10 的指标公式，
\[
  \dim\mathcal M(H)=n+1,
  \qquad H\in\mathcal H_{\mathrm{reg}}.
\]
这是因为，对代表常值映射同伦类 $A_0=0$ 的每个光滑映射
$u:(\mathbb D,\partial\mathbb D)\to(\mathbb C^n,L)$，丛偶
$(\mathbb D\times\mathbb C^n,u^*TL)$ 的\BookRecordIndex{T0062}边界马斯洛夫指标均为零。

现设 $H\in\mathcal H_{\mathrm{reg}}$，并考察\BookRecordIndex{T0283}求值映射
\[
  \operatorname{ev}:\mathcal M(H)\longrightarrow L,
  \qquad \operatorname{ev}(\lambda,u):=u(1,0).
\]
令 $z_0\in L$ 是 $\operatorname{ev}$ 的\BookRecordIndex{T0409}正则值。则集合
\[
  \mathcal M(H;z_0):=\operatorname{ev}^{-1}(z_0)\subset\mathcal M(H)
\]
是光滑紧一维流形，其边界为
\[
  \partial\mathcal M(H;z_0)
  =\{(\lambda,u)\in\mathcal M(H)\mid
      \lambda\in\{0,1\},\ u(1,0)=z_0\}.
\]
由引理~\ref{lem:9-2-4}(i)，\BookRecordIndex{T0246}模空间 $\mathcal M(H)$ 没有 $\lambda=1$ 的元素。因此
$\mathcal M(H;z_0)$ 的边界仅由一点 $(0,u_0)$ 组成，其中
$u_0(s,t)\equiv z_0$。这与紧一维流形的分类矛盾。因此，不存在非常值全纯圆盘的假设必定
错误。定理得证。
\end{proof}

\begin{remark}\label{rem:9-2-5}
\begin{enumerate}[label=\textup{(\roman*)}]
\item 上述论证的一个细化表明，对每个紧嵌入\BookRecordIndex{T0219}拉格朗日子流形
$L\subset\mathbb C^n$，存在\BookRecordIndex{T0242}马斯洛夫数
\[
  \mu(v)\le n+1
\]
的非常值全纯圆盘
$v:(\mathbb D,\partial\mathbb D)\to(\mathbb C^n,L)$。为看出这一点，需要把\BookRecordIndex{T0192}紧性定理
~\ref{thm:5-3-1} 推广到全纯圆盘。关键在于，同伦类在\BookRecordIndex{T0141}格罗莫夫收敛下保持不变（见
Frauenfelder \srcite{120}）。于是，\BookRecordIndex{T0192}紧性的失效会给出一个（连通的）非空有限全纯圆盘
族
\[
  u_j:(\mathbb D,\partial\mathbb D)\longrightarrow(\mathbb C^n,L),
  \qquad j=1,\ldots,N,
\]
以及一个扰动方程的解
$u_0:(\mathbb D,\partial\mathbb D)\to(\mathbb C^n,L)$，使\BookRecordIndex{T0242}马斯洛夫数之和为零：
\[
  \sum_{j=0}^N\mu(u_j)=0.
\]
而对一般\BookRecordIndex{T0157}哈密顿扰动，\BookRecordIndex{T0166}横截性说明当 $\mu<-n-1$ 时，扰动方程的解\BookRecordIndex{T0246}模空间为空。因此
\[
  \sum_{j=1}^N\mu(u_j)=-\mu(u_0)\le n+1,
\]
从而 $\mu(u_1),\ldots,\mu(u_N)$ 中至少一个不超过 $n+1$。所以，如果不存在
$\mu\le n+1$ 的全纯圆盘，\BookRecordIndex{T0246}模空间便是紧的，又会像定理~\ref{thm:9-2-1} 的上述证明那样
导出矛盾。

\item 定理~\ref{thm:9-2-1} 也适用于 Stein 区域中\BookRecordIndex{T0215}可位移的闭\BookRecordIndex{T0219}拉格朗日子流形（Stein
区域是带全局\BookRecordIndex{T0111}多重次调和函数的正合\BookRecordIndex{T0358}辛流形；见下文定义~\ref{def:9-2-6}）。
\end{enumerate}
\end{remark}
\subsection*{凸性}

为了准备下一条关于拉格朗日交点的定理，我们先讨论凸性的概念。Eliashberg 和 Gromov
\srcite{100} 引入了这个概念，以便把 \BookRecordIndex{T0003}$J$-全纯曲线理论推广到非紧\BookRecordIndex{T0358}辛流形。因为在非紧
情形，需要一个假设来保证 \BookRecordIndex{T0003}$J$-全纯曲线不能逃出 $M$ 的某个适当紧子集，从而仍可应用
\BookRecordIndex{T0139}格罗莫夫紧性（见推论~\ref{cor:9-2-11}）。

\begin{definition}\label{def:9-2-6}
称非紧\BookRecordIndex{T0358}辛流形 $(M,\omega)$（在无穷远处）\emph{凸}，如果存在偶对 $(f,J)$，其中 $J$ 是
$\omega$-\BookRecordIndex{T0348}相容\BookRecordIndex{T0088}殆复结构，而 $f:M\to[0,\infty)$ 是恰当光滑函数，并且在 $M$ 的某个紧
子集之外，对每个 $x$ 与每个 $v\in T_xM$ 都有
\begin{equation}\label{eq:9-2-3}
  \omega_f(v,Jv)\ge0,
  \qquad \omega_f:=-d(df\circ J).
\end{equation}
这样的函数 $f$ 称为\emph{\BookRecordIndex{T0111}多重次调和函数}。若 $(M,\omega)$ 是带边界的紧\BookRecordIndex{T0358}辛流形，则称
$M$ 的边界是\emph{凸的}，如果存在偶对 $(f,J)$，其中
$J\in\mathcal J(M,\omega)$，函数 $f:M\to(-\infty,0]$ 在边界附近满足
式~\eqref{eq:9-2-3}，且 $\partial M=f^{-1}(0)$。称这样的\BookRecordIndex{T0088}殆复结构 $J$\emph{适配于边界}。
\end{definition}

\begin{remark}\label{rem:9-2-7}
有些作者对\BookRecordIndex{T0111}多重次调和函数采用更强的定义，要求 $\omega_f=\omega$ 在整个 $M$ 上成立。
有些流形满足弱凸性假设却不满足强凸性假设；见习题~\ref{ex:9-2-14}。定义
~\ref{def:9-2-6} 的一个优点是，两个凸流形的乘积仍然凸。关于辛凸性与复凸性之间极其
密切联系的全面论述，见 Cieliebak--Eliashberg \srcite{63}。
\end{remark}

下一引理给出流形具有凸边界的一个几何判据。

\begin{lemma}\label{lem:9-2-8}
设 $(M,\omega)$ 是带边界的紧\BookRecordIndex{T0358}辛流形，$X\in\operatorname{Vect}(M)$ 是一个在边界附近满足
$\mathcal L_X\omega=\omega$、并在边界处指向外的向量场。则 $(M,\omega)$ 具有凸边界。
\end{lemma}

\begin{proof}
令 $\phi_t\in\operatorname{Diff}(M)$（$t\le0$）为 $X$ 的流。沿边界选取\BookRecordIndex{T0088}殆复结构
$J=\{J(x)\}_{x\in\partial M}$，使得
\begin{equation}\label{eq:9-2-4}
  J(x)X(x)\in T_x\partial M,
  \qquad \omega(v,J(x)X(x))=0,
  \qquad \omega(X(x),J(x)X(x))=1
\end{equation}
对所有 $x\in\partial M$ 与 $v\in T_x\partial M$ 成立。再借助 $X$ 的流把 $J$ 延拓到边界
的一个邻域：
\[
  J(\phi_t(x)):=d\phi_t(x)J(x)d\phi_t(x)^{-1},
  \qquad -\varepsilon<t\le0.
\]
式~\eqref{eq:9-2-4} 的前两个条件说明，相对于度量 $\omega(\,·,J\,·)$，$X$ 垂直于
边界；因此，对每个 $x\in\partial M$，
\sourcefootnote{原书此处把点 $x$ 误写成 $x\in T_xM$；由随后的切空间
$T_xM$ 以及“对每个边界点”的上下文，唯一应为 $x\in\partial M$。}
子空间
$\{v\in T_xM\mid\omega(X(x),v)=0\}$ 在 $J(x)$ 下不变。现选取一个光滑函数
$f:M\to(-\infty,0]$，使得
\[
  f(\phi_t(x))=e^t-1,
  \qquad x\in\partial M,\quad -\varepsilon<t\le0.
\]
\sourcefootnote{原书此处印作 $f(\phi_t(x))=e^t$，但同时规定
$f:M\to(-\infty,0]$ 且 $\partial M=f^{-1}(0)$；在 $t=0$ 时原式给出 $f=1$。唯一\BookRecordIndex{T0348}相容的
归一化是 $e^t-1$。}
则在边界附近，相对于度量 $\omega(\,·,J\,·)$ 有 $\nabla f=X$。于是
$df=-\iota(JX)\omega$，故
\[
  -d(df\circ J)=d\iota(X)\omega=\mathcal L_X\omega=\omega.
\]
引理得证。
\end{proof}

\begin{exercise}\label{ex:9-2-9}
\begin{enumerate}[label=\textup{(\roman*)}]
\item 证明：满足 $\mathcal L_X\omega=\omega$ 的向量场 $X$ 的流 $\phi_t$ 把 $\omega$ 伸缩
$e^t$ 倍，即 $\phi_t^*\omega=e^t\omega$。这样的 $X$ 称为\BookRecordIndex{T0237}刘维尔向量场。

\item 若\BookRecordIndex{T0237}刘维尔向量场 $X$ 定义在 $\partial M$ 附近并与边界横截，证明
$1$-形式 $\alpha:=\iota(X)\omega$ 是 $\partial M$ 上的接触形式。再者，若 $X'$ 是另一个
与 $\partial M$ 横截且指向同一侧的\BookRecordIndex{T0237}刘维尔向量场，则两个接触形式 $\alpha,\alpha'$
同痕，即可由 $\partial M$ 上的一族接触形式相连。因此 $\partial M$ 上诱导出唯一的接触
结构。具有此性质的边界 $\partial M$ 称为\emph{\BookRecordIndex{T0187}接触型}。

\item 证明：若 $Q$ 是 $M$ 中任意闭超曲面，且限制 $\omega|_Q$ 正合，则积分
\[
  \int_Q\alpha\wedge\omega^{n-1}
\]
与满足 $\omega|_Q=d\alpha$ 的 $1$-形式 $\alpha$ 的选取无关。
\end{enumerate}
\end{exercise}

下一引理表明，定义~\ref{def:9-2-6} 中的凸性概念抓住了复情形的重要特征。仍以
$\Delta:=\partial_s^2+\partial_t^2$ 表示标准 Laplace 算子。

\begin{lemma}\label{lem:9-2-10}
设 $J$ 是 $M$ 上的 $\omega$-\BookRecordIndex{T0348}相容\BookRecordIndex{T0088}殆复结构，$\Omega\subset\mathbb C$ 是开集，
$u:\Omega\to M$ 是 \BookRecordIndex{T0003}$J$-全纯曲线，而 $f:M\to\mathbb R$ 是一个在 $u$ 的像的邻域上满足
式~\eqref{eq:9-2-3} 的光滑函数。则 $f\circ u$ 是次调和的：
\[
  \Delta(f\circ u)=\omega_f(\partial_su,J\partial_su)\ge0.
\]
\end{lemma}

\begin{proof}
对 \BookRecordIndex{T0003}$J$-全纯曲线，恒等式
\begin{equation}\label{eq:9-2-5}
  \Delta(f\circ u)=\omega_f(\partial_su,J\partial_su)
\end{equation}
可用两种方法证明。简记 $g:=f\circ u$。由于 $u$ 是 $J$-全纯的，
\[
  u^*(df\circ J)=dg\circ i
  =(\partial_sg\,ds+\partial_tg\,dt)\circ i
  =\partial_tg\,ds-\partial_sg\,dt,
\]
因而
\[
  u^*\omega_f
  =-d\,u^*(df\circ J)
  =-d(\partial_tg\,ds-\partial_sg\,dt)
  =(\Delta g)\,ds\wedge dt.
\]
\sourcefootnote{原书在第二行微分号内把前一行的差
$\partial_tg\,ds-\partial_sg\,dt$ 误印成了和；保留负号时只有这里写成差，才能得到
$(\Delta g)\,ds\wedge dt$。}
这等价于式~\eqref{eq:9-2-5}。

式~\eqref{eq:9-2-5} 的第二种证明基于恒等式
\begin{equation}\label{eq:9-2-6}
  \langle\nabla_\xi\nabla f,\xi\rangle
  +\langle\nabla_{J\xi}\nabla f,J\xi\rangle
  =\omega_f(\xi,J\xi).
\end{equation}
这里 $\nabla$ 是由 $\omega,J$ 确定的度量之 \BookRecordIndex{T0034}Levi--Civita 联络。为证明
式~\eqref{eq:9-2-6}，把恒等式
\[
  d\alpha(X,Y)=\partial_X(\alpha(Y))-\partial_Y(\alpha(X))+\alpha([X,Y])
\]
用于 $1$-形式 $\alpha_f:=-df\circ J=\langle J\nabla f,\,·\rangle$，得到
\[
\begin{aligned}
  \omega_f(X,Y)
  &=\partial_X\langle J\nabla f,Y\rangle
    -\partial_Y\langle J\nabla f,X\rangle
    +\langle J\nabla f,[X,Y]\rangle\\
  &=\langle J\nabla_X\nabla f,Y\rangle
    -\langle J\nabla_Y\nabla f,X\rangle
    +\langle(\nabla_XJ)\nabla f,Y\rangle
    +\langle(\nabla_YJ)X,\nabla f\rangle\\
  &=\langle J\nabla_X\nabla f,Y\rangle
    -\langle J\nabla_Y\nabla f,X\rangle
    -\langle(\nabla_{\nabla f}J)Y,X\rangle
\end{aligned}
\]
对 $X,Y\in\operatorname{Vect}(M)$ 成立；最后一个等式用了式~(C.7.2)。令
$Y=JX$，并利用 $J\nabla_{\nabla f}J$ 是斜自伴随自同态这一事实（比较
式~(C.7.4)），便得到式~\eqref{eq:9-2-6}。其次，每条 \BookRecordIndex{T0003}$J$-全纯曲线
$u:\Omega\to M$ 都是\BookRecordIndex{T0105}调和映射，即满足
$\nabla_s\partial_su+\nabla_t\partial_tu=0$（见注~\ref{rem:4-3-6}）。因此
\[
  \Delta(f\circ u)
  =\langle\nabla_{\partial_su}\nabla f,\partial_su\rangle
   +\langle\nabla_{\partial_tu}\nabla f,\partial_tu\rangle
  =\omega_f(\partial_su,J\partial_su).
\]
最后一个等式来自式~\eqref{eq:9-2-6} 和
$\partial_tu=J\partial_su$。这完成了式~\eqref{eq:9-2-5} 的第二个证明，也证明了引理。
\end{proof}

\begin{corollary}\label{cor:9-2-11}
设 $(M,\omega)$ 是具有凸边界的紧\BookRecordIndex{T0358}辛流形，而 $J$ 是任意适配于边界的 $\omega$-\BookRecordIndex{T0348}相容\BookRecordIndex{T0088}殆复
结构。设 $\Sigma$ 是无边界的连通黎曼曲面，$U\subset M$ 是 $\partial M$ 的邻域，且
$u:\Sigma\to M$ 是光滑映射，其在 $u^{-1}(U)$ 上的限制是 $J$-全纯的。则
\[
  u(\Sigma)\cap\partial M\ne\varnothing
  \quad\Longleftrightarrow\quad
  u(\Sigma)\subset\partial M.
\]
\end{corollary}

\begin{proof}
集合 $\Sigma_0:=\{z\in\Sigma\mid u(z)\in\partial M\}$ 是闭的；由引理
~\ref{lem:9-2-10} 和\BookRecordIndex{T0079}次调和函数的\BookRecordIndex{T0266}平均值不等式，它也是开的。因此它要么为空，要么等于
$\Sigma$。推论得证。
\end{proof}

\begin{example}\label{ex:9-2-12}
$\mathbb C^n$ 上的函数 $f(z)=|z|^2/4$ 是多重次调和的。它的梯度向量场
$X(z)=z/2$ 满足 $\mathcal L_X\omega=\omega$。
\end{example}

\begin{example}\label{ex:9-2-13}
设 $L$ 是紧黎曼流形，以 $\omega_{\mathrm{can}}$ 表示余切丛
$M:=T^*L$ 上的典范辛结构。则纤维中的径向向量场 $X$ 是 $T^*L$ 上的全局\BookRecordIndex{T0237}刘维尔
向量场。为看出这一点，把 $T^*L$ 的元素写成 $(x,y)$，其中
$x\in L$ 且 $y\in T_x^*L$。在局部坐标中，典范 $1$-形式
$\lambda_{\mathrm{can}}\in\Omega^1(T^*L)$ 与典范\BookRecordIndex{T0363}辛形式
$\omega_{\mathrm{can}}=-d\lambda_{\mathrm{can}}\in\Omega^2(T^*L)$ 为
\sourcefootnote{原书把 $\omega_{\mathrm{can}}$ 的类型误印为
$\Omega^1(T^*L)$；它是 $2$-形式，且紧随其后的局部公式也是二次形式。}
\[
  \lambda_{\mathrm{can}}=\sum_i y_i\,dx^i,
  \qquad
  \omega_{\mathrm{can}}=\sum_i dx^i\wedge dy_i,
\]
而 $X:=\sum_i y_i\partial/\partial y_i$。于是
$\iota(X)\omega_{\mathrm{can}}=-\lambda_{\mathrm{can}}$，故
$\mathcal L_X\omega_{\mathrm{can}}=\omega_{\mathrm{can}}$。$L$ 上的度量 $g$ 在 $T^*L$ 上
诱导函数 $\rho_g$，其在 $(x,y)$ 处的值是 $y$ 的长度平方。由于 $X$ 横截于水平集
$\rho_g=1$，我们可如引理~\ref{lem:9-2-8} 在这个水平集上定义 $J$，再沿 $X$ 的流把它
延拓到零截面 $0_L$ 的补集，使 $\rho_g$ 在 $T^*L\setminus0_L$ 上多重次调和。

事实上，度量 $g$ 定义了一个典范的 $J$，使 $\rho_g$ 处处 $J$-多重次调和。这可由如下
公式验证。在局部坐标中，$T^*L$ 上的度量为
\[
\begin{aligned}
  |(\xi,\eta)|^2
  ={}&\sum_{i,j}\xi^ig_{ij}\xi^j\\
  &+\sum_{i,j}
    \left(\eta_i-\sum_{\kappa,\lambda}
      \Gamma^\lambda_{i\kappa}(x)y_\lambda\xi^\kappa\right)
    g^{ij}
    \left(\eta_j-\sum_{k,\ell}
      \Gamma^\ell_{jk}(x)y_\ell\xi^k\right),
\end{aligned}
\]
其中 $\Gamma^\ell_{jk}(x)$ 是 Christoffel 符号。相应的\BookRecordIndex{T0088}殆复结构 $J$ 把 $(x,y)$ 处纤维
的切空间映到其正交补；它由 $J(\xi,\eta)=(\widehat\xi,\widehat\eta)$ 给出，其中
\[
\begin{aligned}
  \widehat\xi^i
  &:=-\sum_jg^{ij}
     \left(\eta_j-\sum_{k,\ell}
       \Gamma^\ell_{jk}(x)y_\ell\xi^k\right),\\
  \widehat\eta_j-
    \sum_{k,\ell}\Gamma^\ell_{jk}(x)y_\ell\widehat\xi^k
  &:=\sum_kg_{jk}\xi^k.
\end{aligned}
\]
于是函数 $f:T^*L\to[0,\infty)$，
\[
  f(x,y):=\frac12|y|^2,
\]
相对于 $J$ 是多重次调和的，因为 $df\circ J=\lambda_{\mathrm{can}}$。
\end{example}

\begin{exercise}\label{ex:9-2-14}
设 $\pi:E\to\Sigma$ 是闭定向二维流形 $\Sigma$ 上次数为 $k$ 的复 Hermitian 线丛，
$P\subset E$ 是单位圆丛。于是 $P$ 是主圆丛，因而带有由向量场 $\partial_\theta$ 生成的
$S^1$-作用。选取总面积为一的面积形式 $\sigma\in\Omega^2(\Sigma)$。则 $P$ 上存在
$S^1$-不变的 $1$-形式 $\alpha$，满足
\[
  \alpha(\partial_\theta)=\frac1{2\pi},
  \qquad d\alpha=-k\pi^*\sigma.
\]
（Bott--Tu \srcite{45, Chapter I, §6} 称之为\emph{\BookRecordIndex{T0289}全局角形式}；可用单位分解把局部解
拼接起来构造它。）以 $r$ 表示 $E$ 上的径向变量。当 $k\le0$ 时，证明
\[
  \omega:=\pi^*\sigma+d(\pi r^2\alpha)
\]
是 $E$ 上的\BookRecordIndex{T0363}辛形式。用引理~\ref{lem:9-2-10} 证明它在\BookRecordIndex{T0344}无穷远处凸。证明：只有当 $k<0$
时，$\omega$ 才满足注~\ref{rem:9-2-7} 中较强的凸性条件。当 $k>0$ 时，证明 $\omega$
不是\BookRecordIndex{T0363}辛形式，但存在有界函数 $\beta:[0,\infty)\to[0,\infty)$，它在 $r=0$ 附近等于
$\pi r^2$，且使
$\omega_\beta:=\pi^*\sigma+d(\beta(r)\alpha)$ 成为\BookRecordIndex{T0363}辛形式。不过，$\omega_\beta$ 在
无穷远处绝不凸。

\emph{提示 1：}证明：当 $k>0$ 时，$(E,\omega_\beta)$ 中的水平集
$Q_c:=\{r=c\}$ 是习题~\ref{ex:9-2-9} 意义下的\BookRecordIndex{T0187}接触型。不过，相应的刘维尔场
$X$ 指向区域 $r<c$，
\sourcefootnote{原书此处写作 $r<r_0$，但上下文只定义了水平集参数 $c$，并未定义
$r_0$；相应刘维尔场指向的正是 $Q_c$ 内侧 $r<c$。}
所以本习题 (iii) 中的积分
$\int_Q\alpha\wedge\omega_\beta$ 是 $c$ 的递减函数，形如
$\operatorname{const}-\operatorname{Vol}(\{r\le c\})$。由此\BookRecordIndex{T0326}推出
$(E,\omega_\beta)$ 不可能在\BookRecordIndex{T0344}无穷远处凸。$k=0$ 时的论证类似。

\emph{提示 2：}$P$ 同构于 $E$ 的\BookRecordIndex{T0067}标架丛，$A:=2\pi i\alpha$ 是一个\BookRecordIndex{T0224}联络形式。证明在下
一习题的记号中，
$\omega=\pi^*\sigma+d(\pi r^2\alpha)=\widetilde\omega$。
\end{exercise}

\begin{exercise}\label{ex:9-2-15}
上述讨论可推广到闭黎曼曲面 $(\Sigma,j,\sigma)$ 上更一般的\BookRecordIndex{T0279}纤维化
\[
  \pi:\widetilde M=P\times_GM\longrightarrow\Sigma.
\]
这里 $G$ 是紧 Lie 群，其 Lie 代数为 $\mathfrak g=\operatorname{Lie}(G)$；
$\pi_P:P\to\Sigma$ 是主 $G$-丛；$(M,\omega)$ 是带哈密顿$G$-作用的\BookRecordIndex{T0358}辛流形。把作用
记作 $(g,x)\mapsto gx$，无穷小作用记作
$\mathfrak g\to\operatorname{Vect}(M):\xi\mapsto X_\xi$。设
$\mu:M\to\mathfrak g^*$ 是该作用的等变\BookRecordIndex{T0204}矩映射，使得
\[
  \omega(X_\xi(x),\widehat x)
  =\langle d\mu(x)\widehat x,\xi\rangle,
  \qquad
  \omega(X_\xi(x),X_\eta(x))
  =\langle\mu(x),[\xi,\eta]\rangle
\]
对 $\widehat x\in T_xM$ 与 $\xi,\eta\in\mathfrak g$ 成立。设
$A\in\mathcal A(P)\subset\Omega^1(P,\mathfrak g)$ 是联络 $1$-形式，其曲率为
$F_A\in\Omega^2(P,\mathfrak g)$。对 $p\in P$，以
$\iota_p:M\to\widetilde M$ 表示纤维的包含，即 $\iota_p(x):=[p,x]$。

\begin{enumerate}[label=\arabic*.]
\item 证明 $2$-形式
$\omega_A:=\omega-d\langle\mu,A\rangle\in\Omega^2(P\times M)$ 下沉为
$\widetilde M$ 上的 $2$-形式
$\widetilde\omega_A\in\Omega^2(\widetilde M)$。证明对
$\widehat x_i\in T_xM$ 与 $v_i\in T_pP$，
\[
\begin{aligned}
  \widetilde\omega_A([v_1,\widehat x_1],[v_2,\widehat x_2])
  ={}&\omega\bigl(\widehat x_1+X_{A_p(v_1)}(x),
                   \widehat x_2+X_{A_p(v_2)}(x)\bigr)\\
     &-\langle\mu(x),F_A(v_1,v_2)\rangle.
\end{aligned}
\]
所以 $(\widetilde M,\widetilde\omega_A)$ 是\BookRecordIndex{T0156}哈密顿纤维化，且对 $p\in P$ 有
$\iota_p^*\widetilde\omega_A=\omega$。

\item 假设对 $x\in M$ 与 $v_i\in T_pP$，
\begin{equation}\label{eq:9-2-7}
  \pi_P^*\sigma(v_1,v_2)>0
  \quad\Longrightarrow\quad
  \langle\mu(x),F_A(v_1,v_2)\rangle\le0.
\end{equation}
证明 $\widetilde\omega:=\pi^*\sigma+\widetilde\omega_A$ 是 $\widetilde M$ 上的\BookRecordIndex{T0363}辛形式。

\item 设 $J$ 是 $M$ 上一个 $G$-不变的 $\omega$-\BookRecordIndex{T0372}驯服\BookRecordIndex{T0088}殆复结构。
\sourcefootnote{原书写作“$X$ 上”的\BookRecordIndex{T0088}殆复结构，但本习题中并未定义流形 $X$；$J$ 的定义域
只能是带 $G$-作用的纤维 $M$。}
证明 $\widetilde M$ 上存在唯一的\BookRecordIndex{T0088}殆复结构 $\widetilde J$，使
$\pi:\widetilde M\to\Sigma$ 和 $\iota_p:M\to\widetilde M$ 对每个 $p\in P$ 都全纯。证明
$\widetilde J$ 是 $\widetilde\omega$-\BookRecordIndex{T0372}驯服的。

\item 设 $K\subset M$ 是 $G$-不变紧集，而 $f:M\to[0,\infty)$ 是 $G$-不变恰当光滑
函数，它在 $M\setminus K$ 上多重次调和，并满足
\begin{equation}\label{eq:9-2-8}
  df(x)JX_{\mu(x)}(x)\ge0,
  \qquad x\in M\setminus K.
\end{equation}
\sourcefootnote{原书式~\eqref{eq:9-2-8} 把 $\mu(x)\in\mathfrak g^*$ 直接作为只对
$\xi\in\mathfrak g$ 定义的向量场 $X_\xi$ 的下标；除非另行指定一个
$\operatorname{Ad}$-不变内积来识别 $\mathfrak g^*\cong\mathfrak g$，该式没有类型意义。
原书未给出这种识别，故这里保留其最低解释度的记号。}
证明函数 $\widetilde f:\widetilde M\to[0,\infty)$，
$\widetilde f([p,x]):=f(x)$，在紧集 $\widetilde K:=P\times_GK$ 的补集上满足
式~\eqref{eq:9-2-3}。（\srcite{135} 用不等式~\eqref{eq:9-2-8} 得到\BookRecordIndex{T0356}辛涡旋方程的
凸性。）
\end{enumerate}
\end{exercise}

\subsection*{拉格朗日交点}

下一条定理断言，\BookRecordIndex{T0219}拉格朗日子流形 $L$ 与它的任意哈密顿变形之间存在交点。它与
Arnold 猜想的一个拉格朗日版本有关。

\begin{theorem}[Gromov]\label{thm:9-2-16}
设 $(M,\omega)$ 是无边界\BookRecordIndex{T0358}辛流形，并假设它在\BookRecordIndex{T0344}无穷远处凸。设 $L\subset M$ 是紧\BookRecordIndex{T0219}拉格朗日
子流形，且 $[\omega]$ 在 $\pi_2(M,L)$ 上消失。若
$\psi:M\to M$ 是\BookRecordIndex{T0161}哈密顿辛微分同胚，则
\[
  \psi(L)\cap L\ne\varnothing.
\]
\end{theorem}

Gromov 在 \srcite{160} 中证明至少存在一个交点。Floer 在 \srcite{113,114} 中证明了
更强的定理：在定理~\ref{thm:9-2-16} 的假设下，只要 $L$ 与 $\psi(L)$ 横截相交，就有
\[
  \#\bigl(L\cap\psi(L)\bigr)\ge\dim H_*(L;\mathbb Z_2).
\]
因此，他在 $\omega$ 于 $\pi_2(M,L)$ 上消失的情形解决了 Arnold 猜想的拉格朗日版本。
Floer 的论文考察的是紧\BookRecordIndex{T0358}辛流形；不过，只要 $(M,\omega)$ 在\BookRecordIndex{T0344}无穷远处凸，证明可逐字移植
到非紧情形。这里给出的论证同样受到 Floer 思想的启发。我们用全纯带直接证明定理
~\ref{thm:9-2-16}，而不通过 Gromov \srcite{160, 2.3.B3'} 的“8 字技巧”从定理
~\ref{thm:9-2-1} \BookRecordIndex{T0326}推出它。

定理~\ref{thm:9-2-16} 中 $[\omega]$ 在 $\pi_2(M,L)$ 上消失的假设不能去掉。例如，若
一条可缩嵌入圆周在二维环面上围成的面积严格小于其补集的面积，便可通过\BookRecordIndex{T0158}哈密顿同痕
把它与自身分开。不过，也有 $[\omega]$ 不在 $\pi_2(M,L)$ 上消失而定理
~\ref{thm:9-2-16} 的结论仍成立的情形。最简单的例子是二维球面上的大圆；另一个例子是
$M=\mathbb{CP}^n$ 且 $L=\mathbb{RP}^n$ 的情形（见 Oh \srcite{303} 和
Abreu--Macarini \srcite{6}）。Fukaya--Oh--Ohta--Ono \srcite{128,129} 发展的拉格朗日
\BookRecordIndex{T0132}弗洛尔同调之 $A_\infty$ 版本，以及 Entov--Polterovich \srcite{106} 的辛刚性理论，都
进一步阐明了这个问题。

回忆定理~\ref{thm:9-2-1} 的断言：$\mathbb C^n$ 的紧\BookRecordIndex{T0219}拉格朗日子流形不能是正合的。这也
可由定理~\ref{thm:9-2-16} \BookRecordIndex{T0326}推出。事实上，若 $L\subset\mathbb C^n$ 是紧\BookRecordIndex{T0219}拉格朗日子流形，
则存在\BookRecordIndex{T0161}哈密顿辛微分同胚 $\psi:\mathbb C^n\to\mathbb C^n$，使
$\psi(L)\cap L=\varnothing$。因此，定理~\ref{thm:9-2-16} 中 $[\omega]$ 在
$\pi_2(\mathbb C^n,L)$ 上消失的假设不可能满足，故 $L$ 不正合。这样，定理
~\ref{thm:9-2-16} 的断言也可理解为拉格朗日嵌入的障碍，而非拉格朗日交点的存在定理。

当 $M$ 紧且 $\omega$ 在 $\pi_2(M)$ 上消失时，定理~\ref{thm:9-1-1} 也可由定理
~\ref{thm:9-2-16} \BookRecordIndex{T0326}推出。事实上，把定理~\ref{thm:9-2-16} 应用于乘积流形
$\widetilde M:=M\times M$、\BookRecordIndex{T0363}辛形式
$\widetilde\omega:=\pi_2^*\omega-\pi_1^*\omega$、对角线
$\widetilde L:=\Delta$，以及由 $H$ 生成的辛微分同胚之图
$\widetilde\psi(x,y):=(x,\psi_H(y))$ 即可。事实上，我们将看到定理
~\ref{thm:9-2-16} 与定理~\ref{thm:9-1-1} 的证明基于同一思路。

\BookRecordIndex{T0235}邻近拉格朗日猜想断言：紧流形之余切丛中的每个\BookRecordIndex{T0401}正合拉格朗日子流形，都与零截面
\BookRecordIndex{T0158}哈密顿同痕。这个猜想仍是公开问题，但利用 \BookRecordIndex{T0021}Fukaya 范畴已经取得了很大进展；例如见
Fukaya--Seidel--Smith \srcite{131}、Abouzaid \srcite{1} 和 Kragh \srcite{217}。
定理~\ref{thm:9-2-16} 的下一推论说明，每个\BookRecordIndex{T0401}正合拉格朗日子流形至少必须与零截面相交。
这一点由 Gromov \srcite{160} 证明；关于它与各纤维相交的第二个断言由
Lalonde--Sikorav \srcite{230} 证明。

\begin{corollary}[Gromov、Lalonde--Sikorav]\label{cor:9-2-17}
设 $L$ 是紧流形，$\Lambda\subset T^*L$ 是\BookRecordIndex{T0401}正合拉格朗日子流形（即
$\lambda_{\mathrm{can}}$ 在 $\Lambda$ 上的限制正合）。则 $\Lambda$ 必与 $T^*L$ 的零
截面以及投影 $\pi:T^*L\to L$ 的每条纤维相交。
\end{corollary}

\begin{proof}
首先，由例~\ref{ex:9-2-13}，余切丛在\BookRecordIndex{T0344}无穷远处凸。其次，由于 $\Lambda$ 正合，
$[\omega_{\mathrm{can}}]$ 在 $\pi_2(T^*L,\Lambda)$ 上消失。因此
$(T^*L,\Lambda)$ 满足定理~\ref{thm:9-2-16} 的假设，从而对每个\BookRecordIndex{T0161}哈密顿辛微分同胚
$\psi:T^*L\to T^*L$，均有
$\psi(\Lambda)\cap\Lambda\ne\varnothing$。

反设 $\Lambda$ 不与零截面相交。考察拉格朗日嵌入
$\iota_t:\Lambda\to T^*L$
\sourcefootnote{原书把靶空间印作本推论中未定义的 $M$；缩放公式及前后的\BookRecordIndex{T0158}哈密顿同痕均在
余切丛 $T^*L$ 中，故唯一应为 $T^*L$。}，
\[
  \iota_t(x,y):=(x,ty),
  \qquad (x,y)\in\Lambda.
\]
下面构造一个满足 $\psi_t(\Lambda)=\iota_t(\Lambda)$ 的\BookRecordIndex{T0158}哈密顿同痕
$\psi_t:T^*L\to T^*L$。但当 $t$ 充分大时
$\iota_t(\Lambda)\cap\Lambda=\varnothing$，这与定理~\ref{thm:9-2-16} 矛盾。

为构造 $\psi_t$，先注意，因为 $\iota_1^*\lambda_{\mathrm{can}}$ 正合，存在紧支撑光滑函数
$H_1:T^*L\to\mathbb R$，它在零截面的一个邻域中消失，并满足
\[
  \iota_1^*\lambda_{\mathrm{can}}+d(H_1\circ\iota_1)=0.
\]
定义
\[
  H_t:T^*L\longrightarrow\mathbb R,
  \qquad H_t(x,y):=tH_1(x,y/t).
\]
于是 $H_t\circ\iota_t=tH_1\circ\iota_1$，恒等式
$\iota_t^*\lambda_{\mathrm{can}}=t\iota_1^*\lambda_{\mathrm{can}}$ 又蕴含
\[
  d(H_t\circ\iota_t)
  =-\iota_t^*\lambda_{\mathrm{can}}
  =\omega_{\mathrm{can}}(\partial_t\iota_t,d\iota_t\,·).
\]
故对每个 $q\in\Lambda$，
\[
  X_{H_t}(\iota_t(q))\in\partial_t\iota_t(q)+\operatorname{im}d\iota_t(q).
\]
所以由 $H_t$ 生成的\BookRecordIndex{T0158}哈密顿同痕 $\psi_t:T^*L\to T^*L$ 满足所需的
$\psi_t(\Lambda)=\iota_t(\Lambda)$。

现设存在 $x_0\in L$，使 $\Lambda$ 不与 $x_0$ 上方的纤维相交。选取 $x_0$ 的开邻域
$U\subset L$，使 $\Lambda\cap\pi^{-1}(U)=\varnothing$；再取光滑函数
$f:L\to\mathbb R$，使 $f$ 的全部\BookRecordIndex{T0233}临界点都属于 $U$。考察\BookRecordIndex{T0158}哈密顿同痕
\[
  \psi_t(x,y):=(x,y+t\,df(x)).
\]
当 $t$ 充分大时，该同痕也会把 $\Lambda$ 与自身分开，而这由定理
~\ref{thm:9-2-16} 不可能发生。推论得证。
\end{proof}

定理~\ref{thm:9-2-16} 的证明几乎逐字照搬上一节定理~\ref{thm:9-1-1} 的证明。相关的
中间结果（见命题~\ref{prop:9-1-4}）涉及单位圆盘
\[
  \mathbb D:=\{s+it\mid s^2+t^2\le1\}\subset\mathbb C
\]
\sourcefootnote{原书把单位圆盘的条件误印为 $t^2+t^2\le1$；圆盘坐标为 $s+it$，故首项
唯一应为 $s^2$。}
上函数 $v:\mathbb D\to M$ 的扰动\BookRecordIndex{T0209}柯西--黎曼方程。它们形如
\begin{equation}\label{eq:9-2-9}
  \partial_sv+X_{F_{s,t}}(v)
  +J(v)\bigl(\partial_tv+X_{G_{s,t}}(v)\bigr)=0,
  \qquad v(\partial\mathbb D)\subset L.
\end{equation}
简记
$K:=F\,ds+G\,dt\in\Omega^1(\mathbb D,C^\infty(M))$，并以
$\|R_K\|$ 表示曲率的 \BookRecordIndex{T0028}Hofer 范数（见注~\ref{rem:8-1-7}）。我们只考察在
$\mathbb D\times M$ 中具有紧支撑并满足
\[
  \zeta\in T_z\partial\mathbb D
  \quad\Longrightarrow\quad K_{z,\zeta}|_L\equiv0
\]
的扰动 $K\in\Omega^1(\mathbb D,C^\infty(M))$。这意味着，相对于由\BookRecordIndex{T0157}哈密顿扰动
$K$ 决定的\BookRecordIndex{T0363}辛形式 $\widetilde\omega_K$，$\widetilde M:=\mathbb D\times M$ 中的
$S^1\times L$ 是\BookRecordIndex{T0219}拉格朗日子流形。以
$\mathcal H(L)\subset\Omega^1(\mathbb D,C^\infty(M))$ 表示这类扰动的集合。由于
$[\omega]$ 在 $\pi_2(M,L)$ 上消失，由引理~\ref{lem:8-2-9}，式~\eqref{eq:9-2-9}
的每个解 $v:\mathbb D\to M$（其中 $K\in\mathcal H(L)$）都满足能量上界
\[
  E_K(v)=\int_{\mathbb D}|\partial_sv+X_F(v)|_J^2\,ds\,dt
  \le\|R_K\|.
\]
以 $\mathcal M(L;K)$ 表示式~\eqref{eq:9-2-9} 的全部解所成空间。

\begin{proposition}\label{prop:9-2-18}
设 $(M,\omega,L)$ 如定理~\ref{thm:9-2-16}，$(f,J)$ 如定义~\ref{def:9-2-6}。则对
每个 $K\in\mathcal H(L)$ 和每个 $w\in\partial\mathbb D$，\BookRecordIndex{T0283}求值映射
\[
  \mathcal M(L;K)\longrightarrow L,
  \qquad v\longmapsto v(w)
\]
都是满射。
\end{proposition}

\begin{proof}
固定常数 $c>0$，使式~\eqref{eq:9-2-3} 在区域 $f\ge c$ 中成立，且 $L$ 包含于集合
$f\le c$。以 $\mathcal H_c(L)$ 表示所有如下\BookRecordIndex{T0157}哈密顿扰动
$K=F_{s,t}\,ds+G_{s,t}\,dt\in\mathcal H(L)$ 的集合：$F_{s,t},G_{s,t}$ 都支撑在紧集
$f\le c$ 中。若 $K\in\mathcal H_c(L)$，则由推论~\ref{cor:9-2-11}，
式~\eqref{eq:9-2-9} 的每个解 $v$ 都取值于紧集
\[
  M^c:=\{x\in M\mid f(x)\le c\},
\]
因为式~\eqref{eq:9-2-3} 在 $M\setminus M^c$ 中成立，且 $L\subset M^c$。其次，我们的
假设保证 $M$ 中没有全纯球面，也没有边界落在 $L$ 上的全纯圆盘。因此，由定理
~\ref{thm:4-6-1}，对每个 $K\in\mathcal H_c(L)$，$\mathcal M(L;K)$ 都是紧的。这里把
定理~\ref{thm:4-6-1} 应用于紧流形
$\widetilde M:=\mathbb D\times M^c$ 和\BookRecordIndex{T0219}拉格朗日子流形
$\widetilde L:=S^1\times L$；$\widetilde M$ 配备由\BookRecordIndex{T0157}哈密顿扰动 $K$ 决定的\BookRecordIndex{T0363}辛形式
$\widetilde\omega_K$ 和\BookRecordIndex{T0088}殆复结构 $\widetilde J_K$。

定理~\ref{thm:8-3-1} 可以显然的方式推广到边值问题~\eqref{eq:9-2-9}，并给出如下结论。
记 $\mathcal H_c^{\mathrm{reg}}(L)\subset\mathcal H_c(L)$ 为所有这样的\BookRecordIndex{T0157}哈密顿扰动
$K\in\mathcal H_c(L)$ 之集合：对每个 $v\in\mathcal M(L;K)$，\BookRecordIndex{T0346}线性化算子
\[
  D_v:
  \Omega^0_{v|_{\partial\mathbb D}^*TL}(\mathbb D,v^*TM)
  \longrightarrow\Omega^{0,1}(\mathbb D,v^*TM)
\]
都是满射。则定理~\ref{thm:8-3-1} 的证明表明，
$\mathcal H_c^{\mathrm{reg}}(L)$ 在 $\mathcal H_c(L)$ 中是\BookRecordIndex{T0300}剩余集。它还表明，对
$K\in\mathcal H_c^{\mathrm{reg}}(L)$，\BookRecordIndex{T0246}模空间 $\mathcal M(L;K)$ 是紧光滑流形，且
\[
  \dim\mathcal M(L;K)=n=\dim L.
\]

现考察\BookRecordIndex{T0283}求值映射
\[
  \operatorname{ev}_K:\mathcal M(L;K)\longrightarrow L,
  \qquad \operatorname{ev}_K(v):=v(w),
  \qquad K\in\mathcal H_c(L).
\]
结合定理~\ref{thm:4-6-1} 与定理~\ref{thm:8-3-3} 可见，对任意两个正则扰动
$K_0,K_1\in\mathcal H_c^{\mathrm{reg}}(L)$，\BookRecordIndex{T0283}求值映射
$\operatorname{ev}_{K_0}$ 与 $\operatorname{ev}_{K_1}$ 由紧配边相连。因此
$\operatorname{ev}_K$ 的模 $2$ 度与 $K\in\mathcal H_c^{\mathrm{reg}}(L)$ 的选择无关。
对零扰动，\BookRecordIndex{T0246}模空间由常值映射组成，故 $\operatorname{ev}_0$ 是微分同胚。因此
$\deg_2(\operatorname{ev}_K)=1$，从而只要
$K\in\mathcal H_c^{\mathrm{reg}}(L)$，$\operatorname{ev}_K$ 就是满射。由于
$\mathcal H_c^{\mathrm{reg}}(L)$ 在 $\mathcal H_c(L)$ 中稠密，显然的\BookRecordIndex{T0192}紧性论证说明，对每个
$K\in\mathcal H_c(L)$，$\operatorname{ev}_K$ 都是满射。又因 $c$ 可任意取大，命题得证。
\end{proof}

\begin{proof}[定理~\ref{thm:9-2-16} 的证明]
固定定义~\ref{def:9-2-6} 中的偶对 $(f,J)$。对每个紧支撑哈密顿函数
\[
  H:\mathbb R\times[0,1]\times M\longrightarrow\mathbb R,
  \qquad(s,t,x)\longmapsto H_{s,t}(x),
\]
\sourcefootnote{原书把该函数的类型误印为
$\mathbb R\times[0,1]\to M$，却仍写三元输入 $(s,t,x)$ 与实值 $H_{s,t}(x)$；其唯一
\BookRecordIndex{T0348}相容的源靶是 $\mathbb R\times[0,1]\times M\to\mathbb R$。}
考察函数 $u:\mathbb R\times[0,1]\to M$ 所满足的方程
\begin{equation}\label{eq:9-2-10}
  \partial_su+J(u)(\partial_tu-X_{H_{s,t}}(u))=0,
  \qquad u(s,0),u(s,1)\in L.
\end{equation}
我们断言：对每个紧支撑哈密顿函数 $H$ 和每个 $x_0\in L$，存在
式~\eqref{eq:9-2-10} 的解 $u:\mathbb R\times[0,1]\to M$，满足 $u(0,0)=x_0$ 和能量
上界
\[
  E(u):=\int_0^1\int_{-\infty}^{\infty}|\partial_su|_J^2\,ds\,dt
  \le\|\partial_sH\|.
\]
为此，考察全纯微分同胚
$\phi:\mathbb R\times[0,1]\to\mathbb D\setminus\{\pm1\}$，
\begin{equation}\label{eq:9-2-11}
  \phi(s,t):=\frac{e^{\pi(s+it)}-i}{e^{\pi(s+it)}+i}.
\end{equation}
不失一般性，可假设对所有 $s,t$，$H_{s,t}$ 的平均值为零；令
$K\in\Omega^1(\mathbb D,C^\infty(M))$ 是唯一满足
$\phi^*K=-H\,dt$ 的\BookRecordIndex{T0157}哈密顿扰动。则
$K\in\mathcal H(L)$ 且 $\|R_K\|=\|\partial_sH\|$。此外，
$v:\mathbb D\to M$ 是式~\eqref{eq:9-2-9} 的解，当且仅当
$u:=v\circ\phi$ 是式~\eqref{eq:9-2-10} 的解。二者能量相同，且
$v(-i)=u(0,0)$。所以断言由命题~\ref{prop:9-2-18} 得证。

其余论证与定理~\ref{thm:9-1-1} 的证明非常相似。考察\BookRecordIndex{T0157}哈密顿扰动
\[
  H_{T;s,t}:=\beta_T(s)H_t:
  \mathbb R\times[0,1]\times M\longrightarrow\mathbb R,
\]
\sourcefootnote{原书在同一页再次漏掉哈密顿函数的流形变量，把定义域印作
$\mathbb R\times[0,1]$；但 $H_t$ 是 $M$ 上的函数，且下文使用其\BookRecordIndex{T0160}哈密顿向量场，故补入
$M$。}
其中 $\beta_T:\mathbb R\to[0,1]$ 是紧支撑光滑\BookRecordIndex{T0188}截断函数，满足
$\beta_T(s)=1$（$|s|\le T$），以及
$\pm\dot\beta_T(s)\le0$（$\pm s\ge0$）。于是，对每个 $T$，存在
式~\eqref{eq:9-2-10} 在 $H=H_T$ 时的解 $u_T$，满足
\[
  E(u_T)\le\|\partial_sH_T\|=2\|H\|.
\]
进而存在形如 $x_T(t)=u_T(s_T,t)$ 的道路 $x_T:[0,1]\to M$，其中
$s_T\in[-T,T]$，并且
\[
  \int_0^1|\dot x_T(t)-X_{H_t}(x_T(t))|^2\,dt
  \le\frac{\|H\|}{T},
  \qquad x_T(0),x_T(1)\in L.
\]
令 $T\to\infty$，得到哈密顿微分方程
$\dot x(t)=X_{H_t}(x(t))$ 的一个解 $x:[0,1]\to M$，它满足边界条件
$x(0),x(1)\in L$。因此 $x(1)=\psi(x(0))$，从而
$x(1)\in L\cap\psi(L)$ 就是所求交点。定理得证。
\end{proof}

\begin{exercise}\label{ex:9-2-19}
在定理~\ref{thm:9-2-16} 的假设下，证明 $L$ 与 $\psi(L)$ 至少存在两个交点。

\emph{提示：}证明：在端点位于 $L$ 且同伦于常值道路的道路空间上，存在良定义的辛作用
泛函。证明，当 $H_{s,t}=H_t$ 时，式~\eqref{eq:9-2-10} 的解正是这个作用泛函相对于由
$J$ 决定的 $L^2$ 内积的梯度流线。
\end{exercise}
